\chapter{Cadre Institutionnel, Stratégique et Organisationnel de l'Office des Changes}

Ce chapitre présente l'organisme d'accueil — l'Office des Changes du Maroc — à travers son historique, ses missions officielles, sa gouvernance, sa stratégie 2025--2029 et le contexte précis dans lequel s'inscrit ce projet de stage.

\section{Présentation générale de l'Office des Changes}

\subsection{Statut juridique et tutelle}

L'Office des Changes est un \textbf{établissement public administratif} doté de la personnalité morale et de l'autonomie financière, placé sous la tutelle du \textbf{Ministère de l'Économie et des Finances}. Il est dirigé par un Directeur nommé par décret royal. Son siège est établi à \textbf{Rabat}, capitale administrative du Royaume du Maroc.

En tant qu'autorité de régulation des opérations de change, l'Office des Changes constitue l'interface institutionnelle entre le système financier marocain et les flux financiers internationaux. Il joue le rôle d'une \textbf{tour de contrôle} de toutes les devises qui traversent les frontières marocaines.

\subsection{Historique et évolution réglementaire}

La réglementation des changes au Maroc a connu plusieurs phases d'évolution, passant d'un régime contraignant à une libéralisation progressive, accordant un intérêt particulier aux attentes et aux besoins réels des opérateurs économiques.

\begin{table}[H]
\centering
\caption{Étapes clés de l'évolution de l'Office des Changes et de la réglementation des changes}
\renewcommand{\arraystretch}{1.4}
\begin{tabular}{|c|p{11cm}|}
\hline
\textbf{Période} & \textbf{Événement} \\
\hline
1944 & Création de l'Office des Changes \\
\hline
Années 1980 & Amorce de la véritable mue de la réglementation — assouplissement progressif \\
\hline
21 jan. 1993 & Adhésion du Maroc à l'article VIII des statuts du FMI relatif à la convertibilité des opérations courantes \\
\hline
1996 & Instauration d'un marché des changes — les exportateurs et importateurs peuvent négocier des taux préférentiels et se couvrir contre le risque de change \\
\hline
2007 & Assouplissement des conditions d'investissement à l'étranger et libéralisation de la gestion des comptes en devises \\
\hline
2010 & Accès libre des citoyens marocains aux devises pour voyages, soins médicaux à l'étranger et commerce électronique \\
\hline
2011 & Codification des règles de change — structuration claire en trois parties : règlements, opérations courantes, opérations en capital \\
\hline
jan. 2020 & Mise en place du Dispositif des Déclarations Bancaires (DDB) et de l'IGOC 2020 \\
\hline
2025 & Lancement de la stratégie 2025--2029 \\
\hline
\end{tabular}
\end{table}

\noindent L'adhésion à l'article VIII du FMI en 1993 a constitué une étape décisive : la convertibilité du dirham adoptée par le Royaume a largement dépassé les exigences minimales de cet article, s'étendant aux investissements étrangers au Maroc, aux financements extérieurs des entreprises marocaines, aux investissements à l'étranger et aux instruments de couverture des risques financiers.

\section{Missions officielles}

Les missions de l'Office des Changes, telles que définies par ses textes fondateurs et publiées sur son site officiel, s'articulent autour de cinq axes :

\subsection{Élaboration de la réglementation des changes}

En étroite collaboration avec le Ministère de l'Économie et des Finances, l'Office participe activement à la mise en œuvre des orientations gouvernementales en matière de change. Il est chargé d'établir les mesures réglementaires et de définir les modalités applicables aux opérations de change, qu'elles concernent les résidents ou les non-résidents. La réglementation est conçue sous forme d'articles et comporte trois parties distinctes : le régime des règlements entre le Maroc et l'étranger, les opérations courantes et les opérations en capital.

\subsection{Contrôle du respect de la réglementation des changes}

L'Office s'assure du respect de la réglementation des changes en vigueur. Il veille à la conformité des opérations par l'examen des documents remis par les opérateurs et les établissements bancaires, ou par des missions d'inspection directe. Il dispose du pouvoir d'accorder des dérogations pour des opérations non couvertes par la réglementation existante et peut constater les infractions et appliquer les mesures répressives prévues par les textes en vigueur.

\subsection{Octroi des agréments de change manuel}

L'Office octroie les agréments de change manuel après vérification des formalités légales requises. Il se garde le droit de contrôler les opérations effectuées par les opérateurs de change manuel, qui sont tenus de mettre à la disposition de ses services de supervision les documents et informations nécessaires sur les opérations effectuées.

\subsection{Établissement des statistiques des échanges extérieurs}

L'Office établit les statistiques des échanges extérieurs et publie mensuellement les \textbf{indicateurs des échanges extérieurs} — balance des paiements, balance commerciale, position financière extérieure globale — qui permettent aux décideurs politiques et économiques de disposer d'un véritable outil d'aide à la prise de décision. Le processus d'établissement des statistiques respecte les délais de production et assure une grande fiabilité des indicateurs, conformément aux normes internationales.

\subsection{Autres missions}

\begin{itemize}
    \item Assumer pleinement son rôle de service public, facilitateur et accompagnateur de ses usagers, en menant une démarche qualité reposant sur la qualité d'accueil, le respect des délais et l'équité de traitement.
    \item Veiller au respect des dispositions de la loi n°\,43-05 relative à la \textbf{lutte contre le blanchiment de capitaux} par les entités soumises à son contrôle.
    \item Dispenser des formations en matière de réglementation des changes au profit des intermédiaires agréés et des associations professionnelles.
    \item Réaliser et publier des études en exploitant les données statistiques des échanges extérieurs.
    \item Échanger les données avec les autres administrations pour partager les informations et assurer un traitement cohérent des questions communes.
\end{itemize}

\section{Gouvernance}

L'Office des Changes s'inscrit dans une dynamique continue de mise en conformité avec les principaux référentiels de bonne gouvernance. Ses instances de gouvernance sont :

\begin{itemize}
    \item \textbf{Le Ministre de l'Économie et des Finances} — autorité de tutelle
    \item \textbf{Le Comité d'Audit} — composé de représentants de la Direction des Entreprises Publiques et de la Privatisation et de la Direction du Trésor et des Finances Extérieures
    \item \textbf{Le Directeur} — \textbf{Driss BENCHIKH}, nommé par décret royal
    \item \textbf{Le Comité de Direction} — coordination stratégique interne
    \item \textbf{Les Comités Spécialisés} : Comité Gestion des Changements, Comité Ressources Humaines, Comité Communication et Relations Extérieures
\end{itemize}

\noindent Le Comité Gestion des Changements valide l'organisation, le planning et l'approche des projets SI, suit leur état d'avancement et prend les décisions de validation technique. C'est dans ce cadre de gouvernance que s'inscrit le projet d'Entity Resolution.

\section{Stratégie 2025--2029}

Élaborée dans le cadre d'une démarche participative, guidée par les Orientations Royales de Sa Majesté le Roi \textbf{Mohammed VI} et par les directives du programme gouvernemental, la stratégie 2025--2029 de l'Office des Changes porte la vision d'être \textit{« une institution au service du développement économique du pays et de la préservation des équilibres extérieurs »}.

Cette stratégie s'articule autour de \textbf{six axes stratégiques} :

\begin{table}[H]
\centering
\caption{Les six axes de la stratégie 2025--2029 de l'Office des Changes}
\renewcommand{\arraystretch}{1.4}
\begin{tabular}{|c|p{4cm}|p{8cm}|}
\hline
\textbf{Axe} & \textbf{Intitulé} & \textbf{Objectif} \\
\hline
1 & Réglementation des changes & Stimuler l'innovation réglementaire pour anticiper les évolutions et renforcer la compétitivité des opérateurs \\
\hline
2 & Contrôle des opérations de change & Renforcer le \textbf{contrôle intelligent} des flux financiers internationaux via des outils technologiques avancés d'analyse et de détection des anomalies \\
\hline
3 & Statistiques des échanges extérieurs & Réinventer la production statistique pour une décision éclairée et une souveraineté informationnelle accrue \\
\hline
4 & \textbf{Digitalisation et gestion de la donnée} & Accélérer la transformation numérique, renforcer la cybersécurité et \textbf{améliorer la qualité des données} \\
\hline
5 & Relation usagers & Optimiser la relation avec les usagers par la simplicité des démarches et la co-construction de services \\
\hline
6 & Gouvernance et développement institutionnel & Instaurer une gouvernance responsable et durable, garantissant la conformité aux normes internationales \\
\hline
\end{tabular}
\end{table}

\noindent Le présent projet de stage s'inscrit directement dans les axes \textbf{2} (contrôle intelligent des flux) et \textbf{4} (digitalisation et qualité des données). En résolvant la fragmentation de l'identité des opérateurs, il contribue à la fois à l'amélioration du contrôle réglementaire et à la fiabilité des données statistiques produites par l'institution.

\section{Structure organisationnelle}

\subsection{Les départements de l'Office des Changes}

L'Office des Changes compte \textbf{sept départements} aux attributions complémentaires :

\begin{itemize}
    \item \textbf{Département Transformation Digitale et Gouvernance de la Donnée (DTDGD)} --- transformation numérique, qualité des données, projets SI, sécurité des systèmes
    \item \textbf{Département Réglementation \& Affaires Juridiques} --- réglementation des changes, dossiers juridiques et contentieux
    \item \textbf{Département Autorisations \& Relations Usagers} --- délivrance des autorisations, contrôle a priori
    \item \textbf{Département Études \& Statistiques} --- balance des paiements, études économiques, publications statistiques
    \item \textbf{Département Supervision} --- contrôle a posteriori sur pièces et missions d'inspection sur place
    \item \textbf{Département Finances, Ressources Humaines \& Moyens Généraux} --- gestion des ressources humaines et financières
    \item \textbf{Département Audit \& Gestion des Risques} --- audit interne et gestion des risques institutionnels
\end{itemize}

\subsection{Le département Transformation digitale et gouvernance de la donnée}

Le \textbf{Département Transformation Digitale et Gouvernance de la Donnée (DTDGD)} constitue le bras technologique et stratégique de l'Office des Changes en matière de numérisation et de valorisation des données. Il est chargé de piloter la transformation digitale de l'institution, d'élaborer et mettre en \oe{}uvre la stratégie data, d'assurer la qualité et la gouvernance des données, et de superviser les projets SI stratégiques.

C'est au sein de ce département, sous la supervision de \textbf{Madame MAJIDI Nadia} (encadrante entreprise) et \textbf{Monsieur MANSOURI Salah-Eddine}, que s'est déroulé ce stage.

\section{Contexte et objectifs du stage}

\subsection{Contexte du projet}

L'Office des Changes reçoit quotidiennement des milliers de déclarations bancaires via le \textbf{Répertoire des Flux Financiers Internationaux (RFF)}, transmises par les établissements de crédit sous format XML conformément à l'Instruction Générale des Opérations de Change (IGOC 2020). Ces déclarations sont stockées dans la table \texttt{RFF.formule\_P\_clean} de la base \texttt{Datamart\_CommerceExterieur}.

Malgré la dématérialisation apportée par le DDB en 2020, une problématique structurelle persiste : \textbf{le même opérateur économique apparaît sous des identités multiples et incohérentes} dans la base de données. Cette fragmentation compromet directement la fiabilité des statistiques produites et l'efficacité des missions de contrôle.

\subsection{Objectifs du stage}

Les objectifs assignés à ce stage de fin d'année sont les suivants :

\begin{enumerate}
    \item \textbf{Analyser la problématique} : étudier la structure de la base \texttt{Datamart\_CommerceExterieur}, identifier les champs affectés par les incohérences d'identifiants, et quantifier l'impact sur la qualité des données.
    \item \textbf{Concevoir un pipeline de nettoyage} : normaliser les champs bruités (\texttt{RC}, \texttt{RAISON\_SOCIALE}, \texttt{DONNEUR\_ORDRE}) avant tout traitement algorithmique.
    \item \textbf{Implémenter un pipeline d'Entity Resolution} : déployer Splink (modèle de Fellegi-Sunter avec algorithme EM) pour regrouper automatiquement les déclarations appartenant au même opérateur.
    \item \textbf{Enrichir par clustering non supervis\'{e} } : compl\'{e}ter le matching probabiliste par du clustering (composantes connexes, Leiden) et des contraintes d'identit\'{e} d\'{e}terministes.
    \item \textbf{Concevoir un mode incrémental} : permettre le traitement des nouvelles déclarations au fil de l'eau sans retraitement de l'historique complet.
    \item \textbf{Évaluer rigoureusement} : construire un jeu de données synthétique avec vérité terrain et mesurer les performances (précision, rappel, F1-score) sur 5 domaines.
\end{enumerate}

\subsection{Compétences mobilisées}

Ce projet mobilise un ensemble de compétences à l'intersection de plusieurs domaines :

\begin{itemize}
    \item \textbf{Ingénierie des données} : modélisation de bases de données SQL Server, nettoyage et normalisation de données, génération de données synthétiques réalistes
    \item \textbf{Machine Learning} : mod\`{e}les probabilistes (Fellegi-Sunter, EM), clustering non supervis\'{e} (Leiden, composantes connexes), \'{e}valuation de mod\`{e}les (pr\'{e}cision, rappel, F1)
    \item \textbf{Connaissance métier} : compréhension du domaine des opérations de change, des formules bancaires RFF, et des contraintes réglementaires de l'OdC
    \item \textbf{Gestion de projet} : documentation technique, architecture logicielle, conception d'un pipeline de production
\end{itemize}

\section{Planning du stage}

Le stage s'est déroulé sur une période de deux mois, du 1\textsuperscript{er} juillet au 31 août 2026, soit 44 jours ouvrables répartis en 8 semaines et demie. Le planning ci-dessous reflète la progression \textbf{réelle} du projet, organisée en \textbf{six phases successives} couvrant l'intégralité du cycle de développement : de la recherche initiale jusqu'au déploiement de la solution complète (pipeline ML, API backend, interface utilisateur).

% ===== PLANNING TABLE =====
\begin{table}[H]
\centering
\caption{Planning réel du stage (1\textsuperscript{er} juillet -- 31 août 2026, 44 jours ouvrables). Les phases Ph.3--6 démarrent en chevauchement avec la fin de Ph.2.}
\renewcommand{\arraystretch}{1.3}
\setlength{\tabcolsep}{3pt}
\begin{tabularx}{\textwidth}{
    |>{\raggedright\arraybackslash}p{0.22\textwidth}
    |>{\centering\arraybackslash}p{0.10\textwidth}
    |*{8}{>{\centering\arraybackslash}X|}}
\hline
\rowcolor{black!80}
{\color{white}\textbf{Phase}} &
{\color{white}\textbf{Période}} &
{\color{white}\textbf{S1}} &
{\color{white}\textbf{S2}} &
{\color{white}\textbf{S3}} &
{\color{white}\textbf{S4}} &
{\color{white}\textbf{S5}} &
{\color{white}\textbf{S6}} &
{\color{white}\textbf{S7}} &
{\color{white}\textbf{S8}} \\
\hline
\rowcolor{teal!25}
\textbf{Ph.1} --- Recherche \& Compréhension &
S1--S2 &
\cellcolor{teal!60} &
\cellcolor{teal!60} &
&
&
&
&
&
\\
\hline
\rowcolor{orange!20}
\textbf{Ph.2} --- Implémentation ML &
S2--S6 &
&
\cellcolor{orange!50} &
\cellcolor{orange!50} &
\cellcolor{orange!50} &
\cellcolor{orange!50} &
\cellcolor{orange!50} &
&
\\
\hline
\rowcolor{blue!15}
\textbf{Ph.3} --- Backend API Flask &
S5--S6 &
&
&
&
&
\cellcolor{blue!45} &
\cellcolor{blue!45} &
&
\\
\hline
\rowcolor{pink!20}
\textbf{Ph.4} --- Frontend Next.js &
S6--S7 &
&
&
&
&
&
\cellcolor{pink!50} &
\cellcolor{pink!50} &
\\
\hline
\rowcolor{gray!15}
\textbf{Ph.5} --- Déploiement Docker &
S7--S8 &
&
&
&
&
&
&
\cellcolor{gray!50} &
\cellcolor{gray!50}
\\
\hline
\rowcolor{yellow!20}
\textbf{Ph.6} --- Rédaction rapport &
S6--S8 &
&
&
&
&
&
\cellcolor{yellow!50} &
\cellcolor{yellow!50} &
\cellcolor{yellow!50}
\\
\hline
\end{tabularx}
\end{table}

% ===== TASKS AND MILESTONES TABLE =====
\begin{table}[H]
\centering
\caption{Principales tâches réalisées et jalons associés}
\renewcommand{\arraystretch}{1.3}
\setlength{\tabcolsep}{4pt}
\begin{tabularx}{\textwidth}{
    |>{\raggedright\arraybackslash}p{0.20\textwidth}
    |>{\raggedright\arraybackslash}X
    |>{\raggedright\arraybackslash}p{0.25\textwidth}|}
\hline
\rowcolor{black!80}
{\color{white}\textbf{Phase}} &
{\color{white}\textbf{Tâches principales}} &
{\color{white}\textbf{Jalon}} \\
\hline
\rowcolor{teal!15}
\textbf{Ph.1} --- Recherche &
Lecture IGOC/DDB, analyse sch\'{e}ma, \'{e}tat de l'art ER/Splink &
Architecture définie \\
\hline
\rowcolor{orange!10}
\textbf{Ph.2} --- ML &
G\'{e}n\'{e}rateur donn\'{e}es, nettoyage RC, Splink v1 (\'{e}chec), Splink v2, DuckDB &
Pipeline ML (P = 72,7\,\%, R = 100\,\%) \\
\hline
\rowcolor{blue!10}
\textbf{Ph.3} --- Backend &
Routes API, job manager async, métriques, tests &
API Flask (port 8083) \\
\hline
\rowcolor{pink!10}
\textbf{Ph.4} --- Frontend &
Dashboard, composants clusters, intégration REST, tests UI &
Interface validée (port 4000) \\
\hline
\rowcolor{gray!10}
\textbf{Ph.5} --- Docker &
Dockerfiles, Docker Compose, tests end-to-end &
Stack déployée \\
\hline
\rowcolor{yellow!15}
\textbf{Ph.6} --- Rapport &
Chapitres 1--4, révision finale, dépôt &
Dépôt du rapport final \\
\hline
\end{tabularx}
\end{table}
% ===== FIN PLANNING TABLE =====

\noindent Les six phases ne sont pas strictement séquentielles : les phases Backend, Frontend et Rédaction démarrent en chevauchement avec la fin de la phase ML, ce qui reflète la réalité d'un projet où développement et documentation avancent en parallèle. Les barres rouges matérialisent l'échec du premier modèle Splink (0,8\,\% de précision) et l'investigation qui a suivi — près d'une semaine de travail, traitée comme information diagnostique plutôt que contournée.

\section*{Conclusion}

Ce chapitre a présenté l'Office des Changes dans toute sa dimension institutionnelle : un établissement public au cœur de la régulation financière marocaine, dont la stratégie 2025--2029 place la digitalisation et le contrôle intelligent comme priorités absolues. Le projet d'Entity Resolution développé dans ce stage s'inscrit directement dans ces axes stratégiques, en apportant une réponse technique concrète à la problématique de qualité des données des opérateurs. Le chapitre suivant analyse en détail cette problématique et la structure de la base de données concernée.
